% ============================================================
% CONCLUSION
% ============================================================
\section{Conclusion and Future Work}

\subsection{Conclusion}

We set out to build an intelligent compiler for medical data-processing programs: a front end that understands the
constraints of the domain, plus a learned component that classifies errors and uses a GAN to cope with scarce labelled
data. The first two parts worked as intended; the third did not.

The compiler for MedScript implements the usual phases---scanner, recursive-descent parser with insertion and panic-mode recovery,
semantic analyser, three-address code, optimiser---and extends the semantic phase with checks that matter in healthcare:
unit-of-measure checking by dimension vectors, physiological plausibility by interval reasoning, taint tracking of protected
health information, flow-sensitive handling of optional values, definite assignment, and division-by-zero detection. All of the
38 unit tests pass, and on 300 generated programs the optimiser preserved observable behaviour on random inputs in every case. Compilation
time grows roughly linearly with program size at about \MsPerKloc{} ms per thousand lines.

For classification, a random forest over 101 compiler-derived features reached \RFAcc\% accuracy on 2{,}000 held-out programs, against
\RuleFirstAcc\% and \RulePrioAcc\% for two rule-based baselines. The gain comes from two situations that rules handle poorly: faults that
produce no diagnostic (inverted clinical rules), and programs in which cascades or several faults make the first diagnostic
misleading. The learning curves show that surface features of the source text are nearly useless on their own (macro-F1 0.44), so
the value of the learner depends on the compiler beneath it.

The GAN result is negative, and we think it is more useful reported than hidden. With 5 to 100 real programs per faulty class, conditional
GANs (BCE and WGAN-GP) gave no improvement over no augmentation for a random forest and a clear loss for an MLP, while Gaussian
jitter and SMOTE gave gains of up to four points of macro-F1. A direct check of the synthetic vectors explains why: they were recognised as the requested class
only about 4\% of the time, and class-balanced training batches raised that only to chance level. A GAN discriminator trained on valid programs did
turn out to be a good unsupervised detector of the faults the compiler misses (AUROC about 0.98), but that is a narrower use than the one the project began with.

\textbf{Limitations.} The corpus is generated, and the test programs come from the same generator as the training programs, so the accuracy figures describe the
generator's faults rather than the faults of real clinical code. The pattern features rely on a hand-written vocabulary. MedScript is a small DSL, not a language
that clinical programmers use today. The optimiser is deliberately simple. These limits are discussed in Section~5.8.

\subsection{Future Work}

\begin{itemize}
    \item \textbf{Real data.} Evaluate on hand-written MedScript programs written by people who did not see the generator, and, more ambitiously, port the
    checks (units, PHI taint, optional values) to a mainstream clinical-scripting language such as Python or R through an AST-based analyser.
    \item \textbf{Better generative models.} Try tabular generators such as CTGAN with proper handling of discrete features, pretrain on the abundant
    valid class, and generate \emph{programs} rather than feature vectors, for instance by training a model over the mutation operators, so that synthetic examples can also be compiled and checked.
    \item \textbf{Learned localisation.} Predict not only the error family but the line at which a soft fault such as an inverted comparison occurs, for example with a
    graph neural network over the syntax tree.
    \item \textbf{Richer semantics.} Add implicit unit conversion with checked conversion factors, interval analysis through loops, inter-procedural taint with function
    summaries, and a three-valued treatment of missing values.
    \item \textbf{A stronger back end.} Global value numbering, propagation across basic blocks, and code generation to a real target such as LLVM IR, so that the pipeline is
    complete from source to machine code.
    \item \textbf{Human evaluation.} Study whether the root-cause estimate shortens the time a developer needs to find a fault compared with raw diagnostics.
\end{itemize}

\newpage
